% =====================================================================
%  H. Brøchner
%  Den nyere Philosophies Historie i Grundrids.
%  (= Philosophiens Historie i Grundrids, 2. Deel.)
%  Kjøbenhavn: P. G. Philipsens Forlag, 1874.   291 pp.
%  Printer: Bianco Lunos Bogtrykkeri.
%
%  DIPLOMATIC TRANSCRIPTION (Danish).
%
%  This is the SECOND and final part of the work whose first part is
%  transcribed at texts/brochner/philosophiens-historie-1 (Den græske
%  Philosophies Historie i Grundrids, 1873, 274 pp.).  The two parts are
%  bound together in the copy scanned and share one PDF, but ANDEN DEEL
%  RESTARTS THE PAGINATION AT 1, which is why it is a separate edition
%  here with its own page map.  Nothing in this file may be checked
%  against the sibling's page numbers.
%
%  THE SOURCE IS NOT A KB SCAN.  Every other Brøchner edition in this repo
%  is transcribed from a Royal Danish Library digitisation; this work could
%  not be.  KB's copy (Den Sorte Diamant, shelfmark 14,-29-8°) is
%  UNDIGITISED and catalogued "Digitalisering eller læsesal".  What is
%  transcribed here is the HARVARD copy -- Widener Library, shelfmark
%  Phil 801.18.3, barcode 3 2044 020 596 763 -- digitised by Google and
%  obtained from HathiTrust (item id hvd.hnqnwx) on 5 September 2026.
%  Should the KB copy be digitised later, this edition should be collated
%  against it.  A KB digitisation is also the obvious way to settle the
%  handful of readings §8 records as buried under the annotator's ink.
%
%  THE SCAN.  597 PDF pages carrying BOTH PARTS bound together: 1. Deel at
%  PDF 1--292, 2. Deel at PDF 293--597.  Page images are 1-bit CCITT G4 at
%  600 ppi, one image per page -- there is no low-resolution RGB layer and
%  no separate mask, unlike the KB scans, so `pdftoppm -r 600` is 1:1 and
%  is the whole of the evidence; anything above it interpolates, which
%  helps a reader and invents nothing.  THE PIXEL DIMENSIONS VARY FROM PAGE
%  TO PAGE, the Google capture having been cropped leaf by leaf; do not
%  take any one of them for the page size, and do not compute a
%  measurement from an assumed one.
%
%  THE COPY IS ANNOTATED, and heavily -- see §8.  It carries a manuscript
%  presentation inscription dated Januar 1874 naming K. Kroman (PDF 6, in
%  the first part's front matter).  1-bit binarisation renders ink the same
%  black as type, so ink and type must be told apart by weight, posture and
%  baseline, never by colour.
%
%  This header is the publication of record for the edition: the page map,
%  the errata, the emphasis and normalisation conventions, and the log of
%  printer's defects.  The per-book harness (pagemap.py, ocr.sh, check.py,
%  joints.py, splice.py, spacing.py, RECON.md, BATCH-AGENT.md,
%  RESUME-NOTES.md) is deliberately untracked, so an outside reader
%  checking this edition against the scan gets this header and note.md and
%  nothing else.
%
% ---------------------------------------------------------------------
%  §1.  THE PAGE MAP -- VERIFIED, UNIFORM, NO STEP
% ---------------------------------------------------------------------
%
%      printed 1--291  =  PDF 300--590,  i.e.  PDF = printed + 299.
%
%  Verified by reading the printed folio off the IMAGE at twenty-four
%  sites: PDF 300 (no folio), 301, 302, 310, 318, 340, 350, 370, 400, 420,
%  450, 470, 498, 499, 520, 530, 531, 550, 563, 564, 580, 589, 590, and
%  591 (blank).  Every one agrees.  No offset step is possible in any
%  case: 590 - 300 = 291 - 1, so the 291 folios exactly fill the 291 PDF
%  pages, leaving no room for the doubly-scanned leaf that stepped the
%  first part's offset from +13 to +15.
%
%  Four pages whose folio the OCR could not read were settled on the
%  image, and all four are set-off or show-through from the facing leaf,
%  not type:  PDF 498 reads 199 (the OCR's "199f" is a comma-shaped
%  smear); PDF 530 reads 231 (its "089" is mirror show-through standing
%  above the folio); PDF 563 reads 264 (its "0264" has a crescent of
%  set-off before the 2).
%
%  THE SIGNATURE TRAP, WHICH THIS WORK SETS TWICE.  Printed p. 1 carries
%  no folio -- it is the display opening of Andet Hovedafsnit -- and a
%  bold "1" stands low on the page.  THAT IS THE SHEET SIGNATURE, NOT A
%  FOLIO: "2" follows on PDF 316, "3" on PDF 332, at the same x, every
%  sixteenth page.  Anyone re-deriving the offset from it corrupts the
%  whole part.  The first part sets exactly the same trap on its own
%  printed p. 1, and its header records it.
%
%  An ink sweep of the folio band over all 291 body pages returns a folio
%  everywhere except PDF 300, so printed 1 is the ONLY page with the folio
%  suppressed.
%
%  FRONT MATTER of Anden Deel, all unnumbered in the print:
%      PDF 293  general title page of the whole work, "2. Deel."
%      PDF 294  part title page
%      PDF 295  imprint (Bianco Lunos Bogtrykkeri)
%      PDF 296--297  Indhold af anden Deel
%      PDF 298--299  Trykfeil og Rettelser -- THREE lists; see §5
%      PDF 300  body begins = printed p. 1
%
%  BACK MATTER: NONE.  Printed 291 carries nine lines, closes with a
%  double rule, and nothing follows -- no Register, no colophon, no
%  ornament, no catchword.  PDF 591--595 carry no type at all (ink counts
%  of 22, 0, 25, 3 and 1 against 12,000--32,000 on a page of text); PDF 596
%  is the Harvard date-due slip and PDF 597 is blank.  Første Deel has a
%  Register; this part has none, and the absence is a fact about the
%  edition, not a gap in the transcription.
%
% ---------------------------------------------------------------------
%  §2.  WHAT THE TWO WITNESSES ARE WORTH, AND WHICH LEADS
% ---------------------------------------------------------------------
%
%  The PDF's text layer is GOOGLE's, not KB's ABBYY, and ITS WORD ORDER IS
%  BROKEN: single words float into the margins and clauses arrive detached
%  from their sentences.  It is an excellent GLYPH witness and a useless
%  STRUCTURE witness.  So `tesseract -l dan` LEADS here for words, lines
%  and paragraphing, and the Google layer adjudicates letters -- the
%  reverse of the arrangement used on the KB books in this repo.
%
%  THE BOOK IS ANTIQUA, NOT FRAKTUR.  Running the Fraktur model on it
%  produces confident nonsense.
%
%  Both witnesses fabricate as well as lose.  In the first part they
%  invented an opening parenthesis on two pages that have none, and they
%  silently normalise every one of the compositor's accent faults.  A
%  reading that exists only in a witness is not evidence.
%
%  THE HATHITRUST/GOOGLE FURNITURE IS BURNT INTO THE PAGE IMAGE -- a grey
%  strip in the left margin and a footer band.  281 of this part's 291
%  pages end in a stray "Google" in the tesseract pass, and spacing.py
%  flags it as letterspaced on every one.  ocr.sh filters it.  It is scan
%  furniture, not type, and none of it belongs in the edition.
%
% ---------------------------------------------------------------------
%  §3.  STRUCTURE AND THE HEAD HIERARCHY
% ---------------------------------------------------------------------
%
%  Upright roman antiqua throughout; no Fraktur, no long s.  The running
%  head is the folio alone, arabic and centred.  The setting is extremely
%  uniform: 33--35 lines to the page at a line pitch that varies by at
%  most one pixel.  Pages materially shorter than that are section
%  boundaries ONLY -- printed 1 (18 lines), 19 (26), 23 (27), 76 (26), 291
%  (9).  A short page anywhere else does not occur, so a transcriber who
%  meets one should suspect he has lost lines rather than that the book
%  has.
%
%  The levels, with the founts measured on the image at 2--3x:
%
%    L0  Andet Hovedafsnit.            bold Egyptian slab-serif, NOT spaced
%    L1  Den nyere Tids Philosophie.   bold condensed SANS-SERIF
%    L2  I n d d e l i n g .           small bold Egyptian, over-inked
%    L2  ... A f d e l i n g .         bold Egyptian, LETTERSPACED
%          its argument                bold condensed serif roman, unspaced
%    L2  ... K a p i t e l .           bold Egyptian, LETTERSPACED
%          its argument                large regular roman, unspaced
%    L3  I. / II. ... Idealisme.       regular roman, LETTERSPACED, centred;
%                                      the roman numeral itself NOT spaced
%    L4  1. Bacon.  etc.               small bold Egyptian, centred, ending
%                                      in a full stop, with no rule of its
%                                      own.  Where the head names a
%                                      philosopher after a colon the NAME
%                                      is letterspaced inside the bold.
%    L2  S l u t n i n g .             bold Egyptian, LETTERSPACED
%
%  THERE IS A SANS-SERIF IN THIS BOOK, and its existence corrects the
%  first part's header, which asserts from the title pages that the work
%  has none anywhere.  In the BODY the sans occurs EXACTLY ONCE: the L1
%  line on printed 1.  Everywhere else a display line is Egyptian or
%  roman.  The Indhold uses it more freely.  Both statements were measured;
%  the first part's claim is hereby withdrawn for this part and flagged for
%  re-examination in that one.
%
%  RULES -- THREE KINDS, AND ONLY THE FIRST IS TEXT.
%    * The DIVISIONAL RULE, a short centred hairline, is part of the text
%      and is set \divrule.  A component sweep finds one on 74 of the 291
%      pages.  THEY ARE OFTEN BROKEN OR UNDER-INKED: the rule above
%      "II. Den kritiske Idealisme" on printed 137 is invisible at a
%      mid-grey threshold and had to be read by eye.  Do not conclude
%      "no rule" from a faint image.
%    * The FOOTNOTE SEPARATOR is left-flush and thicker, and is LaTeX's
%      business, not the transcriber's.  Not transcribed.
%    * The DOUBLE RULE (\dblrule) stands at EXACTLY TWO SITES in the body:
%      the foot of printed 19, closing Første Afdeling, and the foot of
%      printed 291, closing the book.  A sweep over all 291 pages finds no
%      third.
%
%  TWO FACTS ABOUT HEADS, both of which the first part discovered late and
%  at cost, and both of which hold here:
%    1. A HEAD CAN OPEN A PAGE.  Printed 2, 20, 77, 116 and 192 all begin
%       with the head as the first line of type.
%    2. A DIVISIONAL RULE CAN STAND AT THE FOOT OF A PAGE WITH NOTHING
%       AFTER IT.  Printed 115 ends with its rule and the head opens 116;
%       the same at 191 -> 192.  Never infer a head's position from a rule,
%       or a rule's from a head.  Look at the page.
%  Where a numbered section begins mid-page the rule stands ABOVE the head
%  and there is none below it (printed 37, 42, 51, 54, 58, 66, 88, 94, 128,
%  182, 204, 224, 241).
%
% ---------------------------------------------------------------------
%  §4.  CONVENTIONS OF THIS EDITION
% ---------------------------------------------------------------------
%
%  ORTHOGRAPHY EXACTLY AS PRINTED.  Philosophie, Videnskaben, høiere, aa,
%  Charakteer, capitalised nouns.  Nothing is modernised.
%
%  EMPHASIS IS LETTERSPACING, and it is set \emph{}.  The mechanical
%  detector spacing.py runs on this scan and is USABLE but not decisive in
%  either direction: its runs are real but are often display heads, its
%  single-word flags are frequently real (they are fragments of longer
%  spaced phrases), its recall is incomplete, and it flags the scan's
%  "Google" footer on every page.  Every call in this edition was measured
%  by eye on the image, and every page was read for what the tool did not
%  flag.  The first part reached the same verdict after trusting the runs
%  for a while, and withdrew the trust.
%
%  ITALIC IS RESERVED FOR GREEK.  Latin, German and French are set UPRIGHT
%  ROMAN in this book -- printed 216 sets «wirkliches Geschehen» upright,
%  printed 66 sets (Système de la nature) upright.  Nothing that is not
%  Greek is wrapped in \textit{} anywhere in this file.  This differs from
%  many books of the period and was verified on the image, not assumed.
%
%  GREEK is wrapped \gk{} -- the book's Greek fount is a cursive and
%  Libertinus Serif's Greek is upright, so the macro records the fount as a
%  fact about the print rather than as an emphasis, and is greppable.
%  Where a Greek word is ALSO letterspaced, write \emph{\gk{...}}.
%
%  QUOTATION MARKS ARE GUILLEMETS, «...», AND THIS PART OPENS « AND CLOSES
%  » -- the opposite of the 1869 Bidrag, and the same as the first part.
%  There is no „...“ anywhere.  Fifteen pairs were read on the image and no
%  reversed pair was found; printed 200 nests a pair, «...«Ansich»...».
%  Fifteen of some 287 is a sample and not a census, so the quote balance
%  reported by check.py is a live signal: a nonzero balance means a
%  reversed pair or an OCR loss, and it is to be read on the page and
%  logged in §6, never regularised.
%
%  DASHES: one class only, transcribed ---.  The OCR drops and displaces
%  them and is never to be trusted for a dash.
%
%  THE ɔ: MARK (the old Danish id est, U+0254 + colon) is very common in
%  this part -- it appears on some fifty pages -- and BOTH WITNESSES BADLY
%  UNDER-REPORT IT, reading it as "o:", "9:" or "d:".  It is found by
%  reading, not by grepping.
%
%  THE § SORT IS USED -- "(Hegels Encycl. § 384)" on printed 270, read on
%  the image.  The first part concluded its volume had none and was wrong:
%  the Google layer drops the sort and tesseract reads it as an "8", so
%  "the witnesses do not show it" was mistaken for "it is not there".
%
%  FOOTNOTES are plain \footnote{}, the printed mark noted in a % comment
%  at the site.  \thefootnote is redefined below to *) and **); TWO LEVELS
%  ARE ACTUALLY USED (printed 284 and 290 each carry two notes), which is
%  why footmisc[perpage] is not used -- it resets on the TYPESET page, not
%  the printed one.
%
%  PRINTER'S DEFECTS ARE TRANSCRIBED AS PRINTED, logged in a % comment at
%  the site and listed in §6.  The compositor is never silently corrected.
%  Where the printed form is not a word in the language, the sense-reading
%  is preferred unless the glyph POSITIVELY shows otherwise: the burden of
%  proof is on the claim of a defect, because an error in that direction
%  fabricates evidence about the printing house.  An inventory of some 130
%  confident "wrong sorts" built up on the first part had to be withdrawn
%  in its entirety when it was finally calibrated against words the edition
%  already read the other way.  Absence of a feature is not evidence.
%
%  ANYTHING MANUSCRIPT IS NOT TRANSCRIBED.  See §8.
%
% ---------------------------------------------------------------------
%  §5.  THE ERRATA -- THREE LISTS ON TWO LEAVES, NONE APPLIED
% ---------------------------------------------------------------------
%
%  PDF 298--299 carry "Trykfeil og Rettelser", and it is not one list but
%  three.  All three are transcribed in the back matter of this file, and
%  NOT ONE OF THEM IS APPLIED to the text.
%
%    LIST A, headed "1ste Deel." (PDF 298), 30 entries for pp. 28--261 --
%      corrections to FØRSTE DEEL, not to this part.  It is DISJOINT from
%      the "Rettelser." leaf bound at PDF 290 with the first part itself,
%      so Første Deel has 35 errata across two leaves and the sibling
%      edition's header must point here for the other list.
%    LIST B, headed "2den Deel." (foot of PDF 298 running onto PDF 299),
%      17 entries -- the errata for THIS part.  Its sites: printed 6, 12,
%      25, 33, 55, 63, 76, 102, 103, 106, 109 (twice), 155, 177 (twice),
%      187 (twice).
%    LIST C, headed with the title of Brøchner's 1869 "Bidrag til
%      Opfattelsen af Philosophiens historiske Udvikling" (PDF 299), 24
%      entries for pp. 46--280 -- errata TO A DIFFERENT BOOK, bound in
%      here, and belonging to the record of that edition
%      (texts/brochner/philosophiens-udvikling) rather than this one.
%      Do not merge it into list B.
%
%  THE READER OF THIS COPY APPLIED LIST B IN INK, ON THE TYPE.  See §8.
%
% ---------------------------------------------------------------------
%  §6.  PRINTER'S DEFECTS LOGGED, AND THE RUNNING BALANCES
% ---------------------------------------------------------------------
%
%  [To be filled in as the batches report.  Every entry is transcribed as
%  printed and carries a % comment at its site.]
%
%  Balances at the last check:
%    guillemets  0     (« = ») -- see §4; a drift off zero is a signal
%    „...“       0     and there should never be any
%    parentheses 0     -- the first part ran three openers short and one
%                      closer long, all as printed.  Watch for the same.
%
% ---------------------------------------------------------------------
%  §7.  DISAGREEMENTS BETWEEN THE INDHOLD AND THE PRINTED HEADS
% ---------------------------------------------------------------------
%
%  The contents pages and the chapter heads of this book do not agree, and
%  the disagreement is a datum about the printing, not an error to be
%  normalised.  BOTH READINGS ARE SET, the Indhold in the front matter and
%  the head at its page.
%
%  SYSTEMATIC: the Indhold writes "Første Kapitel:" with a colon and runs
%  the argument on; the printed head writes "Første Kapitel." with a full
%  stop and puts the argument on its own line.  The same for Hovedafsnit
%  and Afdeling.  The "Kapitel" line also changes fount between the two --
%  sans in the Indhold, Egyptian on the page.
%
%  SUBSTANTIVE, and there are two:
%    printed 20   Indhold "Dens Udviklingsformer" (plural)
%                 head    "Dens Udviklingsform."  (singular)
%    printed 204  Indhold "4. Reactionen mod Subjectivismen fra
%                          «Realismen»: Herbart"
%                 head    "4. Realismens Reaction mod den subjective
%                          Idealisme: H e r b a r t ."
%                 -- different wording, and the head carries NO GUILLEMETS.
%
%  In the first part, collation of the Indhold against the finished body
%  showed the contents pages siding with the errata leaf against the body
%  at two of five sites, i.e. the contents were set or corrected after the
%  sheets.  Whether the same holds here is an open question for the end of
%  the volume, and is worth the collation: it is a real fact about the
%  printing and it can only be had this way.
%
% ---------------------------------------------------------------------
%  §8.  THE ANNOTATOR, AND THE TEN PLACES WHERE HE CORRECTED THE TYPE
% ---------------------------------------------------------------------
%
%  This copy is annotated more heavily than the first part's: underlining
%  on some twenty-five pages, cursive inside the type area on printed 113
%  and 151, heavy marginal ink on 144, 145, 151 and 216.
%
%  THE DANGEROUS SITES ARE THOSE WHERE THE READER APPLIED ERRATA LIST B IN
%  INK, OVER THE TYPE -- confirmed at printed 6, 12, 25, 33, 55, 76, 106,
%  109, 155 and 177, and suspected at 63, 102, 103 and 187.  AT FOUR OF
%  THEM THE INK OBSCURES THE TYPE: printed 6, 55, 109 and 155.
%
%  THE PRINTED READING STANDS EVERYWHERE.  Where ink buries a letter, the
%  reading is reconstructed and FLAGGED as a reconstruction at the site;
%  it is never quietly supplied, and the annotator's correction is never
%  adopted, however plainly right it is.  Ink is told from type by weight
%  (a pen line fragments under the 1-bit threshold where a type diacritic
%  is solid), by posture (the glosses are cursive, even-hairline and
%  joined), by baseline (ink ignores it, and welds lines into a connected
%  band), and by position (the glosses sit in the interlinear space).
%
%  This is the second copy in this repo where the annotator has done the
%  errata by hand, and it is worth saying why it matters: his corrections
%  agree with list B, so a transcriber who trusted the page without
%  telling ink from type would silently produce a CORRECTED text and
%  believe it diplomatic.
%
% ---------------------------------------------------------------------
%  §9.  OPEN ITEMS
% ---------------------------------------------------------------------
%
%  1. The α) β) γ) enumerators on printed 151 have not been adjudicated --
%     whether they are the book's cursive Greek sort (and so \gk{}) or
%     upright letters used as numbering.  The first part faced the same
%     question, spot-checked its β at 600 dpi, found it cursive, and
%     REMOVED check.py's exemption for such letters.  check.py here has no
%     exemption either, so a bare Greek letter at that site will be
%     reported until someone looks.
%  2. The enumerator series generally.  The first part's turned out to be
%     a genuine mess -- italic and roman alternating inside one section, a
%     "4." where "d." was meant, a comma for a full stop.  Each site here
%     is to be marked ITALIC-ENUM and the whole series adjudicated once,
%     at the end of the volume.  Do not impose a rule page by page.
%  3. The Indhold/body collation described in §7.
%  4. The first part's header asserts that the work uses no sans-serif
%     anywhere.  §3 above falsifies that for this part.  The claim in the
%     sibling's header should be re-examined against its own pages.
%
% =====================================================================

\documentclass[12pt,a4paper]{book}
\usepackage{fontspec}
\usepackage{libertinus}  % engine-aware: loads libertinus-otf under LuaLaTeX
\usepackage{graphicx}    % for \reflectbox (the old Danish »ɔ:« id-est mark)
\usepackage{newunicodechar}
\newunicodechar{ɔ}{\reflectbox{c}}
\usepackage[danish]{babel}
\usepackage[protrusion=true,expansion=true]{microtype}
\usepackage{setspace}
\usepackage[margin=1.2in]{geometry}
\usepackage{enumitem}
\usepackage{amsmath}
\usepackage{fancyhdr}
\setlength{\headheight}{28pt}
\addtolength{\topmargin}{-13.5pt}
\usepackage{hyperref}

\hypersetup{
  pdftitle={Den nyere Philosophies Historie i Grundrids},
  pdfauthor={H. Brøchner},
  hidelinks
}

% TWO footnote levels are actually used in this part -- printed 284 and 290
% each carry two notes (header §4).  NOT \fnsymbol (math mode, gives
% U+2217) and NOT footmisc[perpage], which resets on the TYPESET page and
% would turn a second note on a printed page into a dagger.
\renewcommand{\thefootnote}{%
  \ifcase\value{footnote}\or *)\or **)\or ***)\or \dag)\or \dag\dag)\else
  \arabic{footnote})\fi}

% A short centred divisional rule, as printed above every display head --
% and sometimes with no head after it at all (header §3).  Present on 74
% of the 291 pages, and often broken or under-inked in the scan.
% NOT the footnote separator, which is left-flush and is LaTeX's job.
%
% NB \divrule expands to a center environment, so text following it needs a
% blank line after it to be indented properly.  The first part found this
% late and left a pass over every \divrule as an open item; do it here as
% the batches land, not at the end.
\newcommand{\divrule}{\begin{center}\rule{2cm}{0.4pt}\end{center}}

% The DOUBLE rule: two parallel hairlines, short and centred.  Distinct
% from \divrule, and in this part it stands at EXACTLY TWO SITES in the
% body -- the foot of printed 19 and the foot of printed 291 -- plus the
% title pages.  The 18 pt gap this sets is wider than the print's; tighten
% it when the edition is next reviewed against the page.
\newcommand{\dblrule}{\begin{center}%
  \rule{2cm}{0.4pt}\\[1.2pt]\rule{2cm}{0.4pt}\end{center}}

% The book's Greek fount is a CURSIVE and Libertinus Serif's Greek is
% upright, so every Greek word is wrapped \gk{...}.  One macro rather than
% a bare \textit{}: it is greppable, it records the FOUNT as a fact about
% the print rather than as an emphasis, and a future engine or font change
% is a one-line correction here.  Where a Greek word is ALSO letterspaced,
% write \emph{\gk{...}}; slant and emphasis are independent.
%
% There is very little Greek in this part -- some 40 characters on five
% pages (printed 2, 5, 18, 89, 151) against 8254 characters in the first
% part.  A batch that meets Greek on any other page should suspect it has
% mis-identified something.
\newcommand{\gk}[1]{\textit{#1}}

\pagestyle{fancy}
\fancyhf{}
\fancyhead[LE,RO]{\thepage}
\fancyhead[RE]{\textit{Den nyere Philosophies Historie i Grundrids}}
\fancyhead[LO]{\textit{\leftmark}}

\setstretch{1.3}

\title{\textbf{Den nyere Philosophies Historie}\\[2pt]
  \large i\\[6pt]
  \textbf{\Large Grundrids}}
\author{Dr.\ H.\ Brøchner\\[4pt]
  \small Professor i Philosophien}
\date{Kjøbenhavn: P.\ G.\ Philipsens Forlag, 1874}

\begin{document}
\maketitle
\thispagestyle{empty}
\newpage

% =====================================================================
%  FRONT MATTER
%
%  The front matter stands OUTSIDE the printed numbering and therefore
%  carries no  % --- p. N ---  marker; check.py's arithmetic excludes it.
%  Each leaf is a separate job, done from the image.
%
%  THE 900-SERIES NUMBERS IN THE MARKERS BELOW ARE NOT PAGE NUMBERS.
%  splice.py matches a fragment to its marker by the FIRST--LAST pair in
%  the fragment's filename, so the matter jobs need a pair of their own,
%  and the 900 series was chosen because it cannot collide with a printed
%  page (this part ends at 291) and because check.py reports anything above
%  291 as a bug if it ever appeared as a  % --- p. N ---  marker.  The
%  fragments are .parts/pp900-902.texfrag, pp903-904 and pp905-906.
% =====================================================================

\chapter*{Titelblade}
\markboth{Titelblade}{Titelblade}
% [text to be added: pp. 900--902]
% PDF 293 general title page ("2. Deel."), PDF 294 part title page,
% PDF 295 imprint.  Read at 600 dpi; founts measured, not guessed.
% RECON §2 carries the readings to collate against.

\chapter*{Indhold af anden Deel.}
\markboth{Indhold}{Indhold}
% [text to be added: pp. 903--904]
% PDF 296--297, two unnumbered pages.  Transcribe FROM THE IMAGE, never
% from the contents pages' OCR layer -- the first part's Indhold
% reconstructed from OCR carried two errors.  RECON §3 carries a full
% reading to collate against, including the levels and the leader-dot
% page figures.  Cross-check every argument against its printed head and
% RECORD BOTH where they differ (header §7).

% =====================================================================
%  ANDET HOVEDAFSNIT.  DEN NYERE TIDS PHILOSOPHIE.   printed 1--275
%
%  The body markers below PARTITION printed 1--291 EXACTLY -- contiguous,
%  no overlap, no gap -- and are not to be renegotiated.  Twenty-three
%  batches, cut at section boundaries where the sections allow and cut
%  short over the stretches the annotator has worked on.
% =====================================================================

\chapter*{Andet Hovedafsnit.}
\markboth{Den nyere Tids Philosophie}{Den nyere Tids Philosophie}
% Printed 1 sets this in three lines with a rule between each:
%   "Andet Hovedafsnit."   bold Egyptian, not letterspaced
%   "Den nyere Tids Philosophie."   BOLD CONDENSED SANS -- the only sans
%       line in the whole body (header §3)
%   "Inddeling."   small bold Egyptian, over-inked, letterspaced
% NB it is Inddeling, NOT Indledning: read letter by letter at 600 dpi and
% confirmed by the Indhold.  pagemap.py's docstring had it wrong.
% Printed 1 has NO FOLIO and carries the sheet signature "1" at its foot,
% which is not to be transcribed and is not a page number (header §1).

% [text to be added: pp. 1--12]
% Inddeling (p. 1); Første Afdeling, Overgangsperiodens Philosophie, opens
% printed 2 AT THE TOP OF THE PAGE.  Errata list B has sites at printed 6
% and 12, and the annotator has applied both IN INK ON THE TYPE -- at
% printed 6 the ink OBSCURES type (header §8).  Greek on printed 2
% (κατ᾽ ἐξοχήν) and 5 (τέλος).

% [text to be added: pp. 13--22]
% Close of Første Afdeling at printed 19, which ends with a DOUBLE RULE --
% one of only two in the body -- and is a short page (26 lines) for that
% reason.  Anden Afdeling opens printed 20 at the top of the page, with
% the Første Kapitel head under it: this is the site of the substantive
% Indhold/head disagreement, Udviklingsformer vs Udviklingsform (header
% §7).  Greek on printed 18 (μανία).  Errata site at printed 25 falls in
% the NEXT batch.

% [text to be added: pp. 23--36]
% Andet Kapitel, Den realistiske Udviklingsrække, opens MID-PAGE on
% printed 23 (a short page, 27 lines) with its rule above the head.
% Section 1. Bacon. runs to 37.  Errata sites at printed 25 and 33, both
% applied in ink on the type.  Footnotes on printed 29, 35 and 36.

% [text to be added: pp. 37--50]
% 2. Hobbes. (37) and 3. Locke. (42), each opening mid-page with the rule
% above.  Footnote on printed 39.

% [text to be added: pp. 51--65]
% 4. Berkeley. (51), 5. Hume. (54), 6. Condillac. (58).  Errata sites at
% printed 55 and 63; at printed 55 THE INK OBSCURES TYPE, and 63 is a
% suspected site not yet confirmed (header §8).

% [text to be added: pp. 66--76]
% 7. Materialismen (Système de la nature) -- the parenthesis is set
% UPRIGHT ROMAN, not italic, and the grave on the è is plainly printed.
% Errata site at printed 76.  Printed 76 is a short page (26 lines),
% closing the chapter.

% [text to be added: pp. 77--93]
% Tredie Kapitel, Den idealistiske Udviklingsrække, opens printed 77 AT
% THE TOP OF THE PAGE, followed by the L3 head "I. Den dogmatiske
% Idealisme." and then the run-in head 1. Descartes.  2. Occasionalismen
% (de la Forge, Geulincx, Malebranche) at 88.  Greek on printed 89
% (Γνῶϑι σεαυτόν) -- note the variant sort ϑ, which the book prints and
% which is set as plain θ nowhere in this edition without a note.

% [text to be added: pp. 94--108]
% 3. Spinoza. (94--115).  Errata sites at printed 102, 103 and 106;
% 102 and 103 are suspected rather than confirmed, 106 is confirmed ink
% on type.

% [text to be added: pp. 109--121]
% End of Spinoza.  Printed 115 ends with a DIVISIONAL RULE AND NOTHING
% AFTER IT, and 4. Leibnitz. opens printed 116 at the top of the page
% (header §3).  Errata: printed 109 carries TWO sites and THE INK
% OBSCURES TYPE there.  Footnotes on printed 120 and 121, and THE NOTE
% BEGUN ON 120 RUNS OVER ONTO THE FOOT OF 121, opening with no mark --
% so the foot of 121 must be read even though it opens with ordinary text.

% [text to be added: pp. 122--136]
% End of Leibnitz; 5. Idealismen i det 18de Aarhundrede (128--137).

% [text to be added: pp. 137--149]
% "II. Den kritiske Idealisme." opens printed 137 -- ITS RULE IS SO FAINT
% IT IS INVISIBLE AT A MID-GREY THRESHOLD and must be read by eye; do not
% report "no rule".  1. Kant. runs 137--182, the longest section in the
% volume.  Footnotes on printed 141, 148 and 149.  Heavy marginal ink on
% printed 144 and 145.

% [text to be added: pp. 150--162]
% Kant continued.  Printed 151 carries the α) β) γ) enumerators -- the
% only enumerator series in Greek in this part, NOT YET ADJUDICATED
% (header §9 item 1) -- and cursive ink inside the type area, and heavy
% marginal ink.  Errata site at printed 155, WHERE THE INK OBSCURES TYPE.

% [text to be added: pp. 163--176]
% Kant continued.

% [text to be added: pp. 177--191]
% End of Kant at 182; 2. Reactionen mod Kriticismen fra Troesphilosophien:
% Jacobi (182--191), whose head letterspaces the NAME inside the bold.
% Printed 177 carries TWO errata sites, applied in ink on the type.
% Footnote on printed 190.  Printed 191 ends with a DIVISIONAL RULE AND
% NOTHING AFTER IT.

% [text to be added: pp. 192--203]
% 3. Kriticismens Udvikling til subjectiv Idealisme: Fichte -- the head
% opens printed 192 at the top of the page.  Errata: printed 187 is a
% suspected double site in the PREVIOUS batch; report if it was found.
% Footnote on printed 200.  Printed 200 also carries the volume's one
% NESTED guillemet pair, «...«Ansich»...».

% [text to be added: pp. 204--213]
% 4. Herbart (204--224).  THE PRINTED HEAD DISAGREES SUBSTANTIVELY WITH
% THE INDHOLD here -- "Realismens Reaction mod den subjective Idealisme:
% Herbart", with NO GUILLEMETS, against the Indhold's "Reactionen mod
% Subjectivismen fra «Realismen»: Herbart" (header §7).  Set the head as
% printed and report it.

% [text to be added: pp. 214--223]
% Herbart continued.  Heavy marginal ink on printed 216, which also
% carries «wirkliches Geschehen» -- German, set UPRIGHT ROMAN, not italic
% (header §4).

% [text to be added: pp. 224--232]
% 5. Den speculative Idealisme: a) Schelling (224--241).  The a) belongs
% to a lettered series running in parallel with the arabic numbering;
% mark it ITALIC-ENUM and do not renumber (header §9 item 2).

% [text to be added: pp. 233--240]
% Schelling continued.

% [text to be added: pp. 241--250]
% 6. b) Hegel (241--275).

% [text to be added: pp. 251--262]
% Hegel continued.

% [text to be added: pp. 263--274]
% Hegel continued to 275.  Footnotes on printed 267 and 270.  PRINTED 270
% CARRIES THE VOLUME'S § SORT, "(Hegels Encycl. § 384)" -- both witnesses
% miss it, so it is found by reading (header §4).

% [text to be added: pp. 275--291]
% Slutning (275--291), its head bold Egyptian and letterspaced.  Footnotes
% on printed 278, 284, 285 and 290, and PRINTED 284 AND 290 EACH CARRY TWO
% NOTES -- the only two-note pages in the volume, and the reason
% \thefootnote is defined for two levels.
%
% HOW THE BOOK ENDS.  Printed 291 carries NINE LINES and closes
% "... der ene kan føre til dens Løsning.", then a short centred DOUBLE
% RULE (\dblrule, the second and last in the body).  Below it: nothing.
% No Register, no colophon, no ornament, no catchword, no signature.
% Verified on PDF 590 and 591 (header §1).

% =====================================================================
%  BACK MATTER OF THE SCAN
%
%  The errata leaves are bound after the Indhold in the print, and are set
%  here at the end so that the three lists stand together and can be read
%  against one another.  NONE OF THEM IS APPLIED (header §5).
% =====================================================================

\chapter*{Trykfeil og Rettelser.}
\markboth{Trykfeil og Rettelser}{Trykfeil og Rettelser}
% [text to be added: pp. 905--906]
% PDF 298--299.  THREE LISTS, and they must not be merged:
%   A  "1ste Deel."  30 entries, pp. 28--261 -- corrects FØRSTE DEEL, and
%      is DISJOINT from the Rettelser leaf bound with that part at PDF 290.
%   B  "2den Deel."  17 entries -- corrects THIS part.  Sites: printed
%      6, 12, 25, 33, 55, 63, 76, 102, 103, 106, 109 (twice), 155,
%      177 (twice), 187 (twice).
%   C  headed with the title of the 1869 "Bidrag til Opfattelsen af
%      Philosophiens historiske Udvikling", 24 entries, pp. 46--280 --
%      ERRATA TO A DIFFERENT BOOK, and belonging to that edition's record.
% Two manuscript ink "?" marks on these leaves (PDF 298 at "Side 34",
% PDF 299 at "Side 63"); ink is not transcribed.
% List C's own title is set with „...“ -- the only such marks anywhere in
% the two leaves, everything else being guillemets.  Transcribe as printed.
% RECON §11 carries all three lists to collate against.

\end{document}
